\chapter*{Заключение}
\addcontentsline{toc}{chapter}{Заключение}

В работе предложен фреймворк SA-NAS --- структурно-зависимый поиск
архитектуры для моделей смеси экспертов, связывающий поиск архитектур
отдельных экспертов с кластерной структурой данных. Фреймворк опирается
на графовую нейросетевую суррогатную модель, предсказывающую качество
архитектуры по её описанию и бинарной кластерной маске и обучаемую в
цикле активного обучения с оценкой неопределённости по Monte-Carlo
dropout. Кластерно-зависимый объектив оптимизируется алгоритмом SGEM,
рассматривающим маршрутизацию как латентную переменную; доказана его
сходимость при адаптивно уточняемом суррогате (Теорема~\ref{thm:sgem}).

На контролируемой модельной задаче EM-вариант восстанавливает идеальное
разбиение линейное/кольцевое облако при $K = 2$ экспертах, подтверждая
корректность постановки при надёжном суррогате. На намеренно разнородной
смеси CIFAR-100 / SVHN, чьи истинные домены задают резкий эталон, SGEM
восстанавливает это разбиение на $92\%$ кластеров \emph{без использования
меток источника}: один эксперт специализируется на чистой
CIFAR-подзадаче, другой собирает SVHN-кластеры. При полученном жёстком
кластерном шлюзе достигается точность $0.620$ на тесте --- выше базового
Random-MoE ($0.547$) и $K$ копий одной DARTS-архитектуры ($0.591$) и в
пределах $1.3$ п.п. от оракульного MoE ($0.633$), которому передано
истинное разбиение. Решающую роль здесь сыграл диагностический анализ:
ранний пайплайн, оценивавший экспертов на всей валидации и сэмплировавший
кластерные маски равномерно, никогда не показывал суррогату почти чистые
подмножества и схлопывал маршрутизацию, тогда как оценка по области
эксперта вместе с доменно-агностичным логравномерным сэмплером масок
устраняет это смещение и позволяет кластерно-зависимому объективу
обнаружить структуру данных самостоятельно.

Результаты получены на двухдоменном бенчмарке при $K = 2$ и умеренном
числе сидов; естественным продолжением являются расширение на большее
число доменов, большие $K$ и более широкий набор сидов. Дальнейшие
направления --- масштабирование пространства поиска, перенос фреймворка
на трансформерные MoE-слои, а также усиление теории кластерно-зависимого
объектива, в частности учёт штрафа балансировки нагрузки в гарантии
сходимости.
